\subsection{THETA3-C20: a simple non-CM abelian surface in one ball covering tower}\label{C20}

The elliptic factors in C16 can be replaced by a geometrically
simple abelian surface with no extra endomorphisms. This supplies a
non-CM simple weight-one source of dimension two, rather than a
product of the previously constructed elliptic sources.

\paragraph{Statement.}
Let
\[
 E:\ y^2=x^3-30x-25,\qquad
 \phi(x)=x^4-60x^2-100x-300.
\]
For any nonzero point $P\in E[7](\overline{\Q})$, put $a=x(P)$
and let $H_a$ be the smooth projective curve with affine equation
\[
 H_a:\quad y^2=(x-a)\phi(x).
 \tag{C20a}
\]
Then $H_a$ has genus two and
\[
 \operatorname{End}_{\overline{\Q}}J(H_a)=\Z.
 \tag{C20b}
\]
In particular $J(H_a)$ is geometrically simple and is not CM.
For Schoen's fixed arithmetic ball surface $Y$ in C15--C16, there
is a connected finite \'etale cover $X\to Y$ with
\[
 \deg(X/Y)\leq1764,\qquad
 \operatorname{Alb}(X)\twoheadrightarrow J(H_a).
 \tag{C20c}
\]
The surjection is understood over $\overline{\Q}$ and up to the
harmless choice of origins in the Albanese construction; its
associated $H^1$ inclusion is a rational Hodge morphism.

\paragraph{The division polynomial and its conjugate reductions.}
The four roots of $\phi$ are the nonzero three-torsion
$x$-coordinates, by C16. Set
\begin{align*}
 u&=x^3-30x-25,\\
 b&=x^6-150x^4-500x^3-4500x^2-3000x+22000,\\
 \psi_3&=3\phi,\qquad
 \psi_5=32u^2b-\psi_3^3,\qquad
 \psi_7=\psi_5\psi_3^3-128u^2b^3.
\end{align*}
The elementary division-polynomial recursion
$\psi_{2m+1}=\psi_{m+2}\psi_m^3-
\psi_{m-1}\psi_{m+1}^3$, with $\psi_2=2y$ and
$\psi_4=4yb$, gives these formulas. The exact source is
Silverman, \emph{The Arithmetic of Elliptic Curves},
second edition (2009), Chapter~III, Exercise~3.7,
pp.~105--106, especially parts (b), (d) and (f).
Thus $\psi_7$ has degree
24, leading coefficient seven, and its roots are exactly the
$x$-coordinates of the nonzero seven-torsion points.
One may equally obtain the same equation by the addition law
and the equality $x(4P)=x(3P)$, excluding the smaller-order
denominators; the displayed recursion fixes the normalization.

Here is an explicit finite-field irreducibility certificate.
The monic reduction of $\psi_7$ modulo 19 is
\begin{align*}
 P_{19}(x)={}&x^{24}-9x^{22}-4x^{21}-9x^{20}
       +8x^{19}+6x^{18}-4x^{16}+5x^{15}\\
 &+8x^{14}-2x^{13}-4x^{12}+7x^{11}-7x^{10}
       -9x^8-4x^7+7x^6\\
 &+4x^5+2x^4-4x^3+8x^2-7.
\end{align*}
Successive nineteenth powers modulo this polynomial give
\[
 x^{19^{24}}\equiv x,\qquad
 \gcd(P_{19},x^{19^8}-x)=
 \gcd(P_{19},x^{19^{12}}-x)=1.
 \tag{C20d}
\]
For an explicit check of the large exponents, the ascending
coefficient lists of the two intermediate remainders are
\[
\begin{split}
 r_8={}&(15,13,2,16,7,18,2,6,14,16,7,8,\\
       &\hspace{14mm}10,3,17,11,17,17,6,18,18,1,12,16),\\
 r_{12}={}&(8,18,13,9,13,1,6,13,4,16,14,12,\\
       &\hspace{14mm}14,3,11,3,12,17,5,0,12,6,0,3).
\end{split}
\]
Euclid's algorithm applied to $P_{19}$ and $r_j-x$ gives
the two unit gcds in (C20d). The first equality says that all
irreducible factors have degree dividing 24. Every proper
divisor of 24 divides 8 or 12, so the two gcds exclude all
proper factors. Hence $\psi_7$ is irreducible over $\Q$.
The independent integer-polynomial certificate
\path{documentation/check_C20_independent.py} implements this
recursion, modular powering and Euclidean division with exact
integer arithmetic. It does not use a numerical factorization.

Modulo 37, $\psi_7$ has the simple roots $7,20,35$.
In particular $K=\Q(a)$ has embeddings into $\Q_{37}$ taking
$a$ to Hensel lifts of 7 and 20. These give two places of the
same field $K$, with residue field $\F_{37}$, and two good
reductions of the same curve (C20a). No choice of unrelated
characteristic-zero curves is being compared.

\paragraph{The two point counts and absolute simplicity.}
For the two reductions, direct enumeration gives
\[
\begin{array}{c|r|r|l}
 a\bmod37 & \#H(\F_{37}) & \#H(\F_{37^2})
                    &\text{Frobenius polynomial}\\ \hline
 7 &35&1297&T^4-3T^3-32T^2-111T+1369\\
 20&37&1361&T^4-T^3-4T^2-37T+1369 .
\end{array}
 \tag{C20e}
\]
The polynomials $(x-7)\phi(x)$ and $(x-20)\phi(x)$ are
squarefree modulo 37, so these are smooth genus-two reductions.
The counts include the one point at infinity of the odd-degree
hyperelliptic equation. To specify completely the finite
calculation, use
$\F_{37^2}=\F_{37}[w]/(w^2-2)$; 2 is a nonsquare modulo 37.
For $z=r+sw$ evaluate $(z-a)\phi(z)$ by polynomial arithmetic
and count its square roots in this field. The sum of these
counts over the $37^2$ choices of $(r,s)$, plus one, gives the
second count. An independent implementation uses
$\chi_{\F_{37^2}}(v)=
\chi_{\F_{37}}(\operatorname{Norm}v)$ for $v\ne0$.
The two implementations agree and their full outputs are
included with the paper.

If $A_1=37+1-\#H(\F_{37})$, then
\[
 A_2=\frac{\#H(\F_{37^2})-37^2-1+A_1^2}{2},
 \qquad
 f(T)=T^4-A_1T^3+A_2T^2-37A_1T+37^2.
\]
This derives the two polynomials in (C20e) from the counts,
rather than assuming them. Both are irreducible modulo 5:
in each quartic quotient $T^{5^4}=T$, while
$\gcd(f,T^{5^2}-T)=1$. The respective ascending coefficient
lists of $T^{25}\bmod f$ are $(3,1,4,2)$ and $(1,2,3,2)$.
Thus both quartics are irreducible over $\Q$.
Their middle coefficients $-32$ and $-4$ are units at 37,
so the two Jacobians are ordinary and simple.

We use the exact finite-field criterion of Howe--Zhu,
\emph{On the existence of absolutely simple abelian varieties
of a given dimension over an arbitrary field},
arXiv:math/0002205v1, Theorem~6, p.~5, with proof pp.~5--6.
For $T^4+cT^3+dT^2+37cT+37^2$, its four exceptional
conditions are
$c=0$, $c^2=37+d$, $c^2=2d$ and $c^2=3d-111$.
None holds for $(c,d)=(-3,-32)$ or $(-1,-4)$.
Hence both reductions are absolutely simple.
By the same work, \S2 p.~3 and Proposition~3(2) with
Lemma~4 p.~4, their geometric endomorphism algebras are
exactly the quartic CM fields
\[
 K_1=\Q[T]/(T^4-3T^3-32T^2-111T+1369),\quad
 K_2=\Q[T]/(T^4-T^3-4T^2-37T+1369).
 \tag{C20f}
\]
Indeed absolute simplicity in the ordinary case keeps
$\Q(\pi^r)=\Q(\pi)$ over every finite residue extension.

\paragraph{Excluding every extra geometric endomorphism.}
The real subfield of a quartic ordinary Frobenius field is
generated by $s=\pi+37/\pi$. Its equation is
$s^2+c s+(d-74)=0$ for the coefficients just used.
Consequently
\[
 K_1^+=\Q(\sqrt{433}),\qquad
 K_2^+=\Q(\sqrt{313}).
 \tag{C20g}
\]
The polynomial discriminants, computed by the resultant with
the derivative, are
\[
 \operatorname{disc}(f_1)=
       2^4\,3^3\,37^2\,433^2,\qquad
 \operatorname{disc}(f_2)=
       2^4\,3^3\,11\,37^2\,313^2.
 \tag{C20h}
\]
Both are nonsquares. A quartic CM field containing an
imaginary quadratic subfield would be the compositum of
that field and its real quadratic subfield, hence a
biquadratic Galois extension. A quartic field containing
two different quadratic subfields is likewise
biquadratic. In a biquadratic extension the Galois
permutations of four embeddings are all even, so the
discriminant is a square. The discriminant of a defining
polynomial differs from the field discriminant by a
square. Thus each $K_i$ has exactly one quadratic
subfield, namely the one displayed in (C20g).
The two fields have no common nontrivial subfield up to
$\Q$-isomorphism: degree two is excluded by (C20g), and
an isomorphism in degree four would identify their unique
quadratic subfields.

Put $B=J(H_a)$. Every geometric endomorphism is defined
over some finite extension of $K$. At either of the two
places above 37 it extends to the good abelian scheme and
specializes injectively. The exact source is Conrad,
\emph{Gross--Zagier Revisited}, MSRI Publications~49
(2004), Theorem~2.1(2), printed p.~70.
After any such finite extension the geometric targets
remain the fields $K_1,K_2$, as proved above. Passing to
all geometric endomorphisms therefore embeds the finite
$\Q$-algebra
$\operatorname{End}_{\overline\Q}(B)\otimes\Q$
into each field. It is itself a field: a finite
$\Q$-subalgebra of a number field has no zero divisors,
and multiplication by a nonzero element is a bijection
of its finite-dimensional vector space. Its degree
divides four and is therefore one, two or four.
The last two degrees have just been excluded.
We obtain $\operatorname{End}_{\overline\Q}^0(B)=\Q$.
The integral endomorphism ring is an order in this
field containing $\Z$, hence is $\Z$. This proves
(C20b). A nontrivial abelian subvariety would give a
nontrivial rational idempotent by Poincar\'e
reducibility, so $B$ is geometrically simple.
A CM abelian surface would have a commutative
endomorphism algebra of dimension four, which is
impossible here.

\paragraph{Realization on the fixed ball covering tower.}
The branch points of (C20a) are infinity, the four
$x$-coordinates of $E[3]\setminus\{0\}$ modulo sign,
and $x(P)$ for a seven-torsion point. They are distinct:
an equality of the last coordinate with a three-torsion
coordinate would identify $P$ with a point of order
dividing three, up to sign. All six belong to
$S_{21}=\pi(E[21])$ of C16. The torsion-branch lemma
therefore gives a connected finite \'etale cover
$D\to C$ of degree at most $2\cdot21^2$ and a finite
map $D\to H_a$ of degree at most four. Since
$C\to C_0$ is \'etale of degree two, C15 supplies
a connected ball cover of degree at most
$4\cdot21^2=1764$ with Albanese quotient $J(H_a)$.
This proves (C20c). \hfill$\square$

The graph of the map to $H_a$ and the Abel--Jacobi
correspondence give the asserted rational $H^1$
realization. Thus this finding belongs to the
Jacobian-source and algebraic-correspondence programme
of Grothendieck's cote~44, scans~7--8; the
polarization and correspondence interpretation is
that of cote~16, scans~43--46. It supplies a new
simple source in this particular covering tower.
It does not identify a nonzero exceptional Weil
projection of an image cycle: that remains the
separate extraction problem in C13--C14.
